% Original: PDF strana 21, donji slajd.
\slajdid{02-s042}
\begin{frame}[label=02-s042]{Neodređena stanja i stalni literali}
\setlength{\parskip}{0pt}
% element: 02-s042:e01
\Oznaka{Proizvoljna vrednost: \(b\) ili \(X\)}\par\vspace{1mm}
Ako funkcija u nekom polju može imati bilo koju vrednost, to polje u minimizaciji koristimo kao \textbf{logičku jedinicu ili logičku nulu}, da dobijemo što jednostavniju formu.
\par\vspace{1mm}
Na primer, za zbir proizvoda možemo ga koristiti kao jedinicu, a za proizvod zbirova kao nulu, ili obrnuto --- prema tome šta omogućava povoljnije sažimanje.
\par\vspace{3mm}
% element: 02-s042:e02
Površine objedinjuju proizvode koje je moguće sažeti na prethodno opisani način.
\par\vspace{2mm}
\begin{columns}[T,onlytextwidth]
\begin{column}{.48\textwidth}
\Oznaka{Zbir proizvoda}\par\vspace{1mm}
Ostaju literali koji se \alert{\textbf{ne menjaju}} na površini.
\end{column}
\begin{column}{.48\textwidth}
\Oznaka{Proizvod zbirova}\par\vspace{1mm}
Ostaju \alert{\textbf{komplementi literala}} koji se \alert{\textbf{ne menjaju}} na površini.
\end{column}
\end{columns}
\par\vspace{3mm}
% element: 02-s042:e03
\Rezultat{Promenljive koje se menjaju na površini ne pojavljuju se u odgovarajućem proizvodu ili zbiru.}
\end{frame}
